\documentclass[11pt,reqno]{amsart}

\usepackage[T1]{fontenc}
\usepackage{lmodern}
\usepackage{microtype}
\usepackage{amsmath,amssymb,mathtools}
\usepackage{enumitem}
\usepackage{booktabs}
\usepackage{array}
\usepackage{xcolor}
\usepackage[colorlinks=true,linkcolor=blue!55!black,citecolor=blue!55!black,
  urlcolor=blue!55!black]{hyperref}

\newtheorem{theorem}{Theorem}[section]
\newtheorem{proposition}[theorem]{Proposition}
\newtheorem{lemma}[theorem]{Lemma}
\newtheorem{corollary}[theorem]{Corollary}
\newtheorem{definition}[theorem]{Definition}
\newtheorem{remark}[theorem]{Remark}

\newcommand{\C}{\mathbb C}
\newcommand{\R}{\mathbb R}
\newcommand{\Q}{\mathbb Q}
\newcommand{\Z}{\mathbb Z}
\newcommand{\N}{\mathbb N}
\newcommand{\A}{\mathbb A}
\newcommand{\HC}{\mathrm{HC}}
\newcommand{\HP}{\mathrm{HP}}
\newcommand{\Dom}{\operatorname{Dom}}
\newcommand{\Tr}{\operatorname{Tr}}
\newcommand{\Str}{\operatorname{Str}}
\newcommand{\ord}{\operatorname{ord}}
\newcommand{\im}{\operatorname{Im}}
\newcommand{\re}{\operatorname{Re}}
\newcommand{\clos}{\operatorname{cl}}
\newcommand{\eps}{\varepsilon}
\newcommand{\cS}{\mathcal S}
\newcommand{\cH}{\mathcal H}
\newcommand{\cA}{\mathcal A}
\newcommand{\cC}{\mathcal C}
\newcommand{\cF}{\mathcal F}
\newcommand{\cO}{\mathcal O}
\newcommand{\id}{\mathrm{id}}

\title[Smoothness--detectability conservation]{Smoothness--detectability conservation for Real-corona orientations}

\author{Michael M.}
\address{Independent Researcher}
\email{}

\subjclass[2020]{Primary 46L87, 58B34; Secondary 11M06, 11M26, 19D55, 46L80}
\keywords{cyclic homology, spectral triples, adeles class space, local index formula,
Real structure, automorphic $L$-functions, analytic rank}

\begin{document}

\begin{abstract}
Let $m$ be an isometric involution preserving the domain of a closed unbounded operator $D$.
We prove the exact conservation law
\[
 D=D_{\mathrm{ev}}+D_{\mathrm{odd}},\qquad
 [D,m]=2D_{\mathrm{odd}}m,
\]
and show that $D$ is bounded precisely when both parity components are bounded.  Consequently,
for unbounded $D$, boundedness of $[D,m]$ is equivalent to simultaneous boundedness of the
detecting odd component and unboundedness of the commuting even component.  We apply this law to
the complete Riemann divisor and to the Connes--Consani--Marcolli cyclic module of the adeles class
space.

Four consequences resolve natural problems left by the bounded Real-corona construction.  First,
the canonical complete-divisor realization admits no reflection-compatible local JLO or residue
expansion detecting its nonzero graded orientation: every positive commutator jet vanishes and the
degree-zero term cancels orbitwise.  Second, the exact strong CCM target is not multiplicatively
closed; hence a direct $K_0$ lift is ill-typed.  The nonzero CCM Lefschetz morphisms are Real-even,
whereas the corona character is anti-Real.  The anti-Real scalar shifts are exactly $zL_c$ with
$\bar z=-z$; choosing the positive unit imaginary scalar gives $iL_c$, which squares to $-1$ at
$c=1$.  Third, the corona idempotent is smooth in
the natural dense endpoint-holomorphic crossed product exactly for finite compactified sectors.
Positive-degree self-dual completions have infinite sectors, so smoothness never occurs; this is
proved for the Riemann completion and every nontrivial primitive self-dual Dirichlet completion.
Fourth, in self-dual analytic families the fixed-sector defect variation equals analytic-rank
variation plus twice the off-central variation, while ordinary coordinate $K_0$ sees only support.

Every mathematical assertion introduced in the paper is formalized without additional axioms in
Lean~4 and mathlib; a theorem-by-theorem trace is supplied with the source.
\end{abstract}

\maketitle
\tableofcontents

\section{Introduction}

The adeles class space carries two structures which are individually natural and jointly
in tension.  Scaling produces an unbounded spectral direction, while the functional equation
exchanges the two members of each nonreal divisor orbit.  In the scalar Hardy globalization, the
exchange is represented by an involutive orientation $m=2e-1$, where $e$ is the upper-sector
corona idempotent.  After adjoining functional-equation reflection, $m$ becomes an ordinary
commutator and remains visible only to a graded trace.  These facts motivate four questions:
whether $m$ has a local index formula, whether its $K$-class lifts to the full CCM cyclic module,
when $e$ is smooth in a natural crossed product, and how the fixed real sector compares with
analytic rank.

The central point of this paper is that the four questions are governed by one law.  If $m^2=1$
and $m$ preserves $\Dom D$, then conjugation by $m$ splits $D$ into commuting and anticommuting
parts
\[
 D_{\mathrm{ev}}=\frac{D+mDm}{2},\qquad
 D_{\mathrm{odd}}=\frac{D-mDm}{2}.                 \tag{1.1}
\]
The commutator sees exactly the second part:
\[
 [D,m]=2D_{\mathrm{odd}}m,qquad
 \|[D,m]x\|=2\|D_{\mathrm{odd}}x\|.                \tag{1.2}
\]
What is less tautological, and essential for the applications, is conservation of boundedness:
$D$ is bounded on its domain if and only if both terms in (1.1) are bounded.  Thus a smooth
orientation cannot lose the unbounded spectral direction; it forces that direction into the
commuting sector, where local commutators do not detect it.

This leads to a two-channel principle.  Free functional-equation orbits carry the odd
Clifford/cyclic orientation.  Fixed orbits form a collateral sector satisfying
\[
                 m_0^2=1-e_{\mathrm{fix}}.           \tag{1.3}
\]
The free channel is killed by ordinary crossed-product traces and recovered by a Clifford parity
shift.  The fixed channel is not shifted: its $K_0$ class records support, while its cyclic rank
records multiplicity.  The difference is the exact multiplicity excess.

\subsection{Main results}

The first theorem is purely operator-theoretic.  A partially defined operator means a linear map
$D:\Dom D\to\cH$ on a linear subspace of a normed complex space.  It is bounded on its domain if
there is $C\ge0$ with $\|Dx\|\le C\|x\|$ for all $x\in\Dom D$.

\begin{theorem}[Smoothness--detectability conservation]\label{thm:conservation}
Let $m$ be an isometric complex-linear involution of a normed complex space $\cH$, let $D$ be a
partially defined complex-linear operator, and suppose $m\Dom D\subseteq\Dom D$.  Define
$D_{\mathrm{ev}}$ and $D_{\mathrm{odd}}$ by \textup{(1.1)} on $\Dom D$.  Then
\begin{align*}
 D&=D_{\mathrm{ev}}+D_{\mathrm{odd}},\\
 D_{\mathrm{ev}}m&=mD_{\mathrm{ev}},\qquad
 D_{\mathrm{odd}}m=-mD_{\mathrm{odd}},\tag{1.4}\\
 \|[D,m]x\|&=2\|D_{\mathrm{odd}}x\|,\tag{1.5}\\
 D\text{ is bounded on }\Dom D
 &\Longleftrightarrow
 \begin{gathered}
 D_{\mathrm{ev}}\text{ is bounded on }\Dom D,\\[-1mm]
 D_{\mathrm{odd}}\text{ is bounded on }\Dom D.
 \end{gathered}                                      \tag{1.6}
\end{align*}
If $D$ is unbounded, then
\[
 [D,m]\text{ is bounded}
 \Longleftrightarrow
 \bigl(D_{\mathrm{ev}}\text{ is unbounded and }
       D_{\mathrm{odd}}\text{ is bounded}\bigr).     \tag{1.7}
\]
\end{theorem}

The next result gives the local answer.  The exact class of local formulas is defined in
Section~\ref{sec:p1}.  It includes arbitrary sums of positive iterated-commutator words, hence the
algebraic shape of JLO and residue expansions, together with a summable reflection-invariant
degree-zero spectral term.

\begin{theorem}[Complete-divisor local no-go]\label{thm:p1}
For the multiplicity-exact complete Riemann divisor and its canonical signed diagonal
realization, every reflection-compatible local series evaluated on the Real-corona orientation
is zero.  In particular, no such local series represents the two-dimensional graded orbit charge
$2$.
\end{theorem}

For the full adeles class space, let $C_\bullet^{\mathrm{red}}(\cA)$ be the reduced cyclic module
of the unitized rational adelic crossed product and let
$C_\bullet^{\mathrm{str}}(C_\Q)$ be the CCM strong target: its diagonal normalized trace is
locally constant in the profinite coordinate and belongs, with every derivative, to every
exponentially weighted Schwartz class in the logarithmic radial coordinate.

\begin{theorem}[Full-adeles landing and Clifford comparison]\label{thm:p2}
The actual reduced rational CCM source maps into $C_\bullet^{\mathrm{str}}(C_\Q)$ in every cyclic
degree.  Its degree-zero strong subspace is neither multiplicatively closed nor two-sided in the
ambient trace-norm field algebra.  Therefore the corresponding linear cyclic quotient admits no
zero-absorbing multiplication preserving products of representatives, and a direct $K_0$ lift of
the scalar corona idempotent is not defined.

Every CCM Lefschetz translation $L_c$ is a nonzero Real-even cyclic morphism.  No $L_c$ is
anti-Real.  More exactly, for $z\in\C$,
\[
                 zL_c\text{ is anti-Real}\quad\Longleftrightarrow\quad \bar z=-z.
\]
For the positive unit imaginary scalar, set
\[
                       \mathcal K_c=iL_c             \tag{1.8}
\]
is anti-Real in every cyclic degree, and
\[
                       \mathcal K_1^2=-\id.          \tag{1.9}
\]
\end{theorem}

For smoothness, compactify a self-dual divisor by the centered reciprocal coordinate
$w=(\rho-\tfrac12)^{-1}$.  The functional equation becomes $w\mapsto-w$, and both height ends
accumulate at $0$.  The natural smooth coefficient algebra consists of restrictions to the
compactified divisor of functions holomorphic on a connected neighborhood of the endpoints.

\begin{theorem}[Endpoint smoothness classification]\label{thm:p3}
Let $Z_+$ and $Z_-$ be nonempty disjoint compactified sectors, exchanged by a holomorphic
involution.  Suppose every infinite sector accumulates at its endpoint inside one connected
holomorphic domain.  The upper-sector characteristic idempotent belongs to the endpoint-
holomorphic coefficient algebra, and equivalently lifts to its reflection crossed product, if and
only if both $Z_+$ and $Z_-$ are finite.

In particular, let $\Lambda$ be an entire self-dual completion with zeros in $0<\re s<1$, a
subquadratic exponential envelope, and logarithmic growth on the positive real ray exceeding every
linear slope.  Then the compactified Real-corona idempotent is not endpoint-holomorphically smooth.
This applies unconditionally to the Riemann completion and to every nontrivial primitive self-dual
Dirichlet completion.
\end{theorem}

Finally, let $F_t$ be a self-dual analytic family with center $c$, and count real zeros in a fixed
reflection-invariant compact window with analytic multiplicity.  Let $d(t)$ be total fixed-sector
multiplicity, $r(t)=\ord_c F_t$ analytic rank, and $u(t)$ total multiplicity of real zeros strictly
above $c$.

\begin{theorem}[Fixed-sector variation]\label{thm:p4}
For all parameters $t_0,t_1$,
\[
 d(t_1)-d(t_0)=r(t_1)-r(t_0)+2\bigl(u(t_1)-u(t_0)\bigr). \tag{1.10}
\]
The central coordinate evaluation of the fixed Real-corona $K_0$ class is
$\mathbf1_{r(t)>0}$.  Hence
\[
 r(t_1)-r(t_0)
 =\mathbf1_{r(t_1)>0}-\mathbf1_{r(t_0)>0}
  +\bigl(\mu(r(t_1))-\mu(r(t_0))\bigr),              \tag{1.11}
\]
where $\mu(n)=n-\mathbf1_{n>0}$ is central multiplicity excess.
\end{theorem}

\subsection{Relation to prior work and scope}

Connes--Consani--Marcolli formulate the Weil explicit formula as a Lefschetz trace on cyclic
homology of a cokernel in the abelian category of cyclic modules \cite{CCM}.  Meyer constructs the
corresponding idele-class representation in a nuclear bornological setting \cite{Meyer}.  Local
index formulas and JLO characters provide the standard route from spectral triples to cyclic
cohomology \cite{CM,JLO,Moscovici}.  Our contribution is an obstruction and replacement theorem:
the complete-divisor orientation cannot be simultaneously locally detecting and smooth in the
intrinsic realization; the exact CCM target is cyclic-linear rather than multiplicative; and the
parity-corrected Lefschetz morphism is forced to be Clifford.

The paper does not prove the Riemann hypothesis, Weil positivity, convergence of finite zeta
spectral triples, or a higher-rank automorphic instance not encoded by the stated analytic
hypotheses.  The smooth algebra in Theorem~\ref{thm:p3} is the endpoint-holomorphic dense algebra
defined in Section~\ref{sec:p3}; no claim is made for every possible smooth completion.

\section{The parity conservation law}\label{sec:law}

Let $V$ be a normed complex vector space.  Let $D:\Dom D\to V$ be complex-linear and let
$m:V\to V$ be complex-linear with
\[
 m^2=1,\qquad \|mx\|=\|x\|,\qquad m\Dom D\subseteq\Dom D. \tag{2.1}
\]
The second application of domain invariance and $m^2=1$ shows $m\Dom D=\Dom D$.

\begin{definition}
On $\Dom D$ define
\[
 D_{\mathrm{ev}}x=\frac12(Dx+mDmx),\qquad
 D_{\mathrm{odd}}x=\frac12(Dx-mDmx).                \tag{2.2}
\]
\end{definition}

\begin{lemma}[Parity identities]\label{lem:parity}
The identities \textup{(1.4)} hold, and
\[
 [D,m]x=2D_{\mathrm{odd}}(mx).                       \tag{2.3}
\]
\end{lemma}

\begin{proof}
Adding the two formulas in (2.2) gives $D$.  Using $m^2=1$ gives
\[
 D_{\mathrm{ev}}mx=\tfrac12(Dmx+mDx)=mD_{\mathrm{ev}}x
\]
and similarly $D_{\mathrm{odd}}mx=-mD_{\mathrm{odd}}x$.  Finally,
\[
 2D_{\mathrm{odd}}mx=Dmx-mDm(mx)=Dmx-mDx=[D,m]x.
\]
\end{proof}

\begin{lemma}[Norm identity]\label{lem:norm}
For every $x\in\Dom D$,
\[
              \|[D,m]x\|=2\|D_{\mathrm{odd}}x\|.
\]
Consequently $[D,m]$ is bounded on $\Dom D$ if and only if $D_{\mathrm{odd}}$ is bounded there.
\end{lemma}

\begin{proof}
Equation (2.3), the odd parity relation, and isometry of $m$ give
\[
 \|[D,m]x\|=2\|D_{\mathrm{odd}}mx\|
 =2\|mD_{\mathrm{odd}}x\|=2\|D_{\mathrm{odd}}x\|.
\]
The boundedness equivalence follows, with constants changed by a factor of two.
\end{proof}

\begin{lemma}[Conservation of boundedness]\label{lem:bounded}
$D$ is bounded on $\Dom D$ if and only if both $D_{\mathrm{ev}}$ and $D_{\mathrm{odd}}$ are
bounded on $\Dom D$.
\end{lemma}

\begin{proof}
If $\|Dx\|\le C\|x\|$, then domain invariance and isometry give
\[
 \|Dmx\|\le C\|mx\|=C\|x\|,
 \qquad \|mDmx\|=\|Dmx\|.
\]
The triangle inequality in (2.2) bounds each parity component by $C\|x\|$; the formal proof uses
the harmless common bound $2C$.  Conversely, if the two components have bounds $C_{\mathrm{ev}}$
and $C_{\mathrm{odd}}$, then (1.4) bounds $D$ by their sum.
\end{proof}

\begin{proof}[Proof of Theorem~\ref{thm:conservation}]
Lemmas~\ref{lem:parity}--\ref{lem:bounded} prove (1.4)--(1.6).  If $D$ is unbounded and $[D,m]$
is bounded, Lemma~\ref{lem:norm} makes $D_{\mathrm{odd}}$ bounded.  Lemma~\ref{lem:bounded} then
forces $D_{\mathrm{ev}}$ to be unbounded.  The converse in (1.7) again follows from
Lemma~\ref{lem:norm}.
\end{proof}

\begin{remark}
Theorem~\ref{thm:conservation} is independent of self-adjointness and compact resolvent.  Those
hypotheses become relevant when the law is applied to spectral triples; the obstruction itself is
already present at the invariant-domain level.
\end{remark}

\section{The complete divisor and local cyclic blindness}\label{sec:p1}

We now specialize to the complete nontrivial divisor of the Riemann completion.  The formal model
keeps every zero with its analytic multiplicity.  Let $I$ be this multiplicity-copy index set,
let $r:I\to I$ be functional-equation reflection, and let
\[
 z_i=\rho_i-\tfrac12,\qquad z_{r(i)}=-z_i.            \tag{3.1}
\]
There are no fixed copies in the nontrivial Riemann divisor.  Choose an orbit orientation
$\eps:I\to\{\pm1\}$ satisfying $\eps(r(i))=-\eps(i)$, and put
\[
 \cH=\ell^2(I),\qquad
 (ma)_i=\eps(i)a_i.                                  \tag{3.2}
\]
The maximal real diagonal operators relevant here are
\[
 (D_{\mathrm{abs}}a)_i=|z_i|a_i,qquad
 (D_{\mathrm{sgn}}a)_i=\eps(i)|z_i|a_i.              \tag{3.3}
\]
Both are self-adjoint and genuinely unbounded.  Reflection $R(a)_i=a_{r(i)}$ commutes with
$D_{\mathrm{abs}}$ and anticommutes with $D_{\mathrm{sgn}}$.  The diagonal orientation $m$
commutes with both.

\begin{proposition}[Extremal complete-divisor realizations]\label{prop:extremal}
Reflection is an all-orders smooth element for $D_{\mathrm{abs}}$.  It has genuinely unbounded
first commutator for $D_{\mathrm{sgn}}$, so no all-orders smooth representative of the reflection
generator exists there.  In contrast, $m$ is all-orders smooth for $D_{\mathrm{sgn}}$, but all its
positive iterated commutators vanish.  Moreover
\[
                        RmR=-m.                      \tag{3.4}
\]
\end{proposition}

\begin{proof}
The coordinate identities following (3.3) preserve the maximal domains because reflection and
$m$ preserve coordinate norms.  Exact commutation promotes the bounded operators to smooth
elements with zero higher commutators.  For the signed operator,
$D_{\mathrm{sgn}}R=-RD_{\mathrm{sgn}}$; hence
\[
 [D_{\mathrm{sgn}},R]=2D_{\mathrm{sgn}}R.
\]
Theorem~\ref{thm:conservation} identifies its norm with twice the unbounded diagonal norm.  Finally
(3.4) is $\eps(r(i))=-\eps(i)$ coordinatewise.
\end{proof}

We make the local-formula statement exact.  For a smooth represented element $a$, write
$\delta^0(a)=a$ and let $\delta^k(a)$, $k>0$, denote any bounded iterated commutator extension used
by the smooth calculus.  A \emph{positive local word} is a finite product
\[
                 \delta^{k_1}(a)\cdots\delta^{k_j}(a) \tag{3.5}
\]
with some $k_l>0$.  A positive local series is an arbitrary indexed sum of scalar multiples of
additive functionals applied to words (3.5).  No convergence hypothesis is needed for the
vanishing result because every summand is zero.  A complete-divisor local series additionally has
a degree-zero term
\[
        \sum_{i\in I}\eps(i)h(i),                    \tag{3.6}
\]
where the summand is absolutely summable and $h(r(i))=h(i)$.

\begin{lemma}[Positive-word extinction]\label{lem:positive}
Every positive local word and every positive local series for $m$ relative to
$D_{\mathrm{sgn}}$ is zero.
\end{lemma}

\begin{proof}
Exact commutation gives $\delta^k(m)=0$ for every $k>0$.  Every positive word contains one such
factor.  Therefore every summand in the series is zero.
\end{proof}

\begin{lemma}[Orbit cancellation]\label{lem:orbitcancel}
If the family in \textup{(3.6)} is summable and $h\circ r=h$, then (3.6) is zero.
The same conclusion holds for every finite reflection-invariant packet without a summability
hypothesis.
\end{lemma}

\begin{proof}
Reindexing the sum by $r$ and using $\eps\circ r=-\eps$ gives $S=-S$, hence $S=0$.  On a finite
packet the same proof is a fixed-point-free involution pairing.
\end{proof}

\begin{proof}[Proof of Theorem~\ref{thm:p1}]
Every complete local series is the sum of (3.6) and a positive local series.  The two terms vanish
by Lemmas~\ref{lem:positive} and \ref{lem:orbitcancel}.  The graded two-orbit representation has
grading matrix $\mathrm{diag}(1,-1)$ and its supertrace on $m$ equals $2$.  Therefore no complete
local series has the required value.
\end{proof}

\begin{remark}
The theorem is a no-go for the canonical, multiplicity-exact, reflection-compatible complete-
divisor realization.  It does not exclude a different spectral triple whose Dirac operator has a
nonzero bounded odd component.  Theorem~\ref{thm:conservation} states the necessary price: all
locally detected orientation data must enter through that bounded odd component.
\end{remark}

\section{The rational CCM cyclic module}\label{sec:ccm}

We recall only the structures needed for the statement.  Let
\[
       \cA=\cS(\A_\Q)\rtimes\Q^\times                \tag{4.1}
\]
denote the rational adelic Bruhat crossed product, with its minimal unitization $\cA^+$.  There are
two cyclic endpoint morphisms $\eps_0^\bullet,\eps_1^\bullet$ induced respectively by zero-adele
evaluation and the product of Lebesgue and normalized finite-adelic Haar integration.  The reduced
source is the augmentation kernel inside
\[
  \ker\eps_0^\bullet\cap\ker\eps_1^\bullet
       \subset C_\bullet(\cA^+).                     \tag{4.2}
\]
This is the cyclic module used in the CCM construction, rather than the full unitized cyclic bar.

At a nonzero adelic base, the Morita representation maps (4.1) to the minimal unitization of the
trace-class algebra on $\ell^2(\Q^\times)$.  Varying the base gives a continuous trace-norm field.
For a cyclic tensor $x_0\otimes\cdots\otimes x_n$, take the normalized trace of the ordered
pointwise product and restrict to the main diagonal.  In the logarithmic radial coordinate $t$,
the strong scalar target requires
\[
  \sup_t e^{\beta t}\left|\frac{d^k}{dt^k}f(t)\right|<\infty
  \quad\text{for every }\beta\in\R\text{ and }k\ge0, \tag{4.3}
\]
together with local constancy in the profinite-unit coordinate.  The inverse image of this
condition is a cyclic submodule, denoted $C_\bullet^{\mathrm{str}}(C_\Q)$.

The analytic input to the landing theorem is the following exact periodic statement.  Our Fourier
transform uses the self-dual additive Haar normalization.

\begin{lemma}[Periodic strong-decay theorem]\label{lem:periodic}
Let $M\ge1$ and let $g_r\in\cS(\R)$ for $r\in\Z/M\Z$.  Put
\[
 \Theta_g(t)=\sum_{r\bmod M}\ \sum_{\substack{n\in\Z\\n\equiv r\bmod M}}
                         g_r(e^t n).                 \tag{4.4}
\]
If
\[
                    g_0(0)=0,\qquad
                    \sum_{r\bmod M}\widehat g_r(0)=0, \tag{4.5}
\]
then $\Theta_g$ is smooth and, for every $\beta\in\R$ and $k\ge0$,
\[
             \sup_{t\in\R}e^{\beta t}|\Theta_g^{(k)}(t)|<\infty. \tag{4.6}
\]
Only the two global cancellations in \textup{(4.5)} are required; the individual Fourier zero
modes need not vanish.
\end{lemma}

\begin{proof}
Write $E=x\,d/dx$.  Termwise differentiation is justified by a uniform Schwartz majorant on a
neighborhood of each $t$, and gives
\[
 \Theta_g^{(k)}(t)=\sum_{r}\sum_{n\equiv r(M)}(E^kg_r)(e^tn). \tag{4.7}
\]
At $t\to+\infty$, the only stationary lattice term is $g_0(0)$; it vanishes by (4.5), while an
arbitrarily high Schwartz seminorm bounds all $n\ne0$.  This proves every positive exponential
weight at that end.

For the opposite end, periodic Poisson summation gives, first for a rank-one family
$g_r=w_r g$,
\[
 \sum_n w_{n\bmod M}g(e^tn)
 =\frac{e^{-t}}{M}\sum_{m\in\Z}\widehat w_m\,widehat g(e^{-t}m/M). \tag{4.8}
\]
Conjugating $E$ through (4.8) replaces it by $A=-(1+E)$ on the Fourier side.  Hence every
derivative is again a weighted lattice sum of a Schwartz function.  Its stationary term vanishes
when $(\sum_rw_r)\widehat g(0)=0$, and Schwartz seminorms give arbitrary negative-end weights.

For a general family, if $\widehat g_0(0)\ne0$ choose $r_0=0$; otherwise choose any index
$r_0$ with $\widehat g_{r_0}(0)\ne0$, if one exists.  Set
\[
 c_r=\frac{\widehat g_r(0)}{\widehat g_{r_0}(0)},\qquad
 h_r=g_r-c_rg_{r_0}.                                \tag{4.9}
\]
Every $h_r$ has Fourier zero mode zero, while $\sum_rc_r=0$ by (4.5).  The stated choice also gives
$h_0(0)=0$: either $h_0=0$, or $c_0=0$ and (4.5) applies.  Thus the $h$-part is a sum
of pointwise-cancelled traces and the rank-one correction satisfies the rank-one hypothesis.
The case in which every Fourier zero mode vanishes is immediate.  This proves (4.6).
\end{proof}

We next record the finite-ad\`elic step in a form which retains cancellation between tensor
summands.  For an algebraic Bruhat tensor
$x\in\cS(\R)\otimes\cS(\A_{\Q,f})$, a profinite unit $b$, and $q\in\Q$, let
$\Phi_x(a,b;q)$ denote evaluation at the rational translate $(aq,qb)$.

\begin{lemma}[Common-period adelic packet]\label{lem:packet}
For every $x$ and $b$ there are $M\ge1$ and a family
$g_r\in\cS(\R)$, $r\in\Z/M\Z$, such that
\begin{align}
 \sum_{q\in\Q}\Phi_x(e^t,b;q)&=\Theta_g(t),                 \tag{4.10}\\
 g_0(0)&=\eps_0(x),                                         \tag{4.11}\\
 \sum_{r\bmod M}\widehat g_r(0)&=M\eps_1(x).               \tag{4.12}
\end{align}
If $\eps_0(x)=0$, the sum in \textup{(4.10)} may be reindexed by $\Q^\times$.
Moreover $b\mapsto(t\mapsto\sum_{q\in\Q^\times}\Phi_x(e^t,b;q))$ is locally constant as a map
to the full function space $\C^\R$.
\end{lemma}

\begin{proof}
For $x=f\otimes\phi$, compact support and local constancy of $\phi$ give a positive denominator
$D$ and an open subgroup $M\widehat\Z$ on which all relevant translates are constant.  Exact Haar
quadrature on
\[
                         \widehat\Z/M\widehat\Z\simeq\Z/M\Z \tag{4.13}
\]
produces residue weights $w_r$ with
\[
 \sum_rw_r=M D^{-1}\!\int_{\A_{\Q,f}}\phi,
 \qquad
 g_r=w_r\,f(D^{-1}\,\cdot).                         \tag{4.14}
\]
The rational orbit decomposes into these residue classes, giving (4.10).  Evaluation at the zero
orbit gives (4.11), and Fourier scaling together with (4.14) gives (4.12), including the normalized
finite-ad\`elic Haar factor $|q|_\infty^{-1}$.

An arbitrary $x$ is a finite sum of pure tensors.  Replace all periods by their product.  Period
refinement repeats each residue statistic by the exact refinement factor, so both sides of
(4.11)--(4.12) remain additive.  This is why endpoint cancellation is imposed only after the
summands have been assembled.  If (4.11) is zero, the $q=0$ term vanishes identically and the
injective inclusion $\Q^\times\hookrightarrow\Q$ gives the stated reindexing.

Finally, the compact fractional-integral support of the finitely many finite-ad\`elic factors has
one open profinite-unit stabilizer.  Right multiplication by that stabilizer fixes every orbit
coefficient simultaneously for all $t$; hence the whole radial profile is locally constant.
\end{proof}

The last reduction is operator-theoretic rather than asymptotic.

\begin{lemma}[Morita trace and cyclic-degree reduction]\label{lem:morita}
For $a\in\ell^1(\Q^\times)$ and $k\in\Q^\times$, let
\[
                 (W_{a,k}\xi)(q)=a(q)\xi(k^{-1}q).   \tag{4.15}
\]
Then $W_{a,k}$ is trace class and
\[
     \|W_{a,k}\|_1=\sum_q|a(q)|,\qquad
     \Tr(W_{a,k})=\mathbf1_{k=1}\sum_q a(q).         \tag{4.16}
\]
The resulting Morita representation of $\cA^+$ is a continuous trace-norm field and is
multiplicative.  If $x=x_0\otimes\cdots\otimes x_n$, its strong diagonal trace and both cyclic
endpoints are respectively
\begin{align}
 \Tr^\natural_n(\rho^{\otimes n+1}x)&=\Tr^\natural_0(\rho(x_0\cdots x_n)),\tag{4.17}\\
 \eps_j^n(x)&=\eps_j(x_0\cdots x_n),\qquad j=0,1.    \tag{4.18}
\end{align}
\end{lemma}

\begin{proof}
The upper bound in (4.16) factors $W_{a,k}$ as a product of two Hilbert--Schmidt weighted shifts
with coefficients $|a|^{1/2}$ and the phase times $|a|^{1/2}$.  For the reverse inequality, pair
the operator against the two orthonormal bases $(\delta_q)_q$ and $(\delta_{kq})_q$ in the trace-
norm dual estimate; the resulting sum is exactly $\sum_q|a(q)|$.  The diagonal matrix coefficient
is zero unless $k=1$, proving the trace formula.  Orbit summability of Bruhat functions puts each
crossed-product monomial in this $\ell^1$ model, and the identity
\[
 W_{a,k}W_{d,l}=W_{q\mapsto a(q)d(k^{-1}q),kl}        \tag{4.19}
\]
is the crossed-product multiplication law.  The exact norm identity plus dominated convergence
gives trace-norm continuity in the class-space base.

Equations (4.17)--(4.18) follow from multiplicativity: ordered tensor multiplication commutes with
the representation and with both algebra characters.  Taking normalized trace yields (4.17).
Linearity extends the argument from pure tensors to every cyclic chain.
\end{proof}

\subsection{The exact strong landing}

\begin{proposition}[All-degree source landing]\label{prop:landing}
The class-space restriction of every chain in the reduced source \textup{(4.2)} belongs to
$C_\bullet^{\mathrm{str}}(C_\Q)$ in every degree.
\end{proposition}

\begin{proof}
Let $X$ be a reduced chain and let $y$ be its ordered total product.  The augmentation-kernel
condition gives zero scalar part for $y$, while (4.18) gives
$\eps_0(y)=\eps_1(y)=0$.  Let $x$ be the identity-scaling Bruhat coefficient of $y$.  The exact
unitized trace formula (4.16) identifies the class-space trace of $y$ with
\[
              (b,t)\longmapsto
              \sum_{q\in\Q^\times}\Phi_x(e^t,b;q).   \tag{4.20}
\]
By Lemma~\ref{lem:packet}, this is $\Theta_g$ for a common-period packet satisfying (4.5).
Lemma~\ref{lem:periodic} proves every estimate (4.3), and the last assertion of
Lemma~\ref{lem:packet} supplies local constancy in $b$.  Thus the degree-zero trace belongs to the
literal strong class-space target.  Equations (4.17)--(4.18) now reduce every cyclic degree to this
same $y$, proving the proposition.
\end{proof}

\subsection{Multiplication fracture}

Let $F$ be the degree-zero ambient algebra of continuous trace-norm fields and let
$S\subset F$ be the exact strong degree-zero subspace.  The next two matrix units live in the
actual Schatten-$1$ Morita fibre.  In the $\Q^\times$ basis, write
\[
                  E_+=E_{1,-1},\qquad E_-=E_{-1,1}.  \tag{4.21}
\]
They are implemented by $\ell^1$-weighted shifts, hence are trace class.

\begin{proposition}[Strong multiplication fracture]\label{prop:fracture}
The constant fields $E_+$ and $E_-$ belong to $S$, but $E_+E_-$ does not.  More precisely,
\[
 \Tr(E_+)=\Tr(E_-)=0,
 \qquad \Tr(E_+E_-)=1.                               \tag{4.22}
\]
Consequently $S$ is neither multiplicatively closed nor two-sided.  There is no multiplication
on $F/S$ which absorbs zero on both sides and sends classes of representatives to the class of
their ambient product.
\end{proposition}

\begin{proof}
The weighted-shift trace is supported on identity scaling.  Each operator in (4.21) has scaling
$-1$, so its trace is zero.  Their product has identity scaling and coefficient sum one, proving
(4.22).  A constant scalar function satisfies all two-ended bounds (4.3) if and only if it is zero:
using any positive exponential weight at $+\infty$ forces a nonzero constant out of the target.
Thus the first two constant fields lie in $S$, while their product does not.

If $S$ were two-sided, the product of its two displayed members would belong to $S$, a
contradiction.  If multiplication descended to $F/S$ with the stated properties, the zero classes
of $E_+$ and $E_-$ would have zero product, while representative multiplication would give the
nonzero class of $E_+E_-$.
\end{proof}

\subsection{Real parity and the Lefschetz correspondence}

The exact strong cyclic module has an even conjugate-linear Real structure $\sigma$.  A cyclic
morphism $T$ is Real-even if $\sigma T=T\sigma$ and anti-Real if
$\sigma T=-T\sigma$.  CCM class-space translation by $c\in C_\Q$ gives an invertible Real-even
cyclic morphism $L_c$.

\begin{lemma}[Parity intersection]\label{lem:parityintersection}
A cyclic morphism which is both Real-even and anti-Real is zero in every degree.
\end{lemma}

\begin{proof}
For every chain $x$, the two parity identities give
$\sigma T x=T\sigma x=-T\sigma x$.  Since $\sigma$ is invertible and $2$ is invertible in $\C$,
$Tx=0$.
\end{proof}

\begin{proposition}[Clifford Lefschetz transgression]\label{prop:clifford}
Every $L_c$ is nonzero and is not anti-Real.  For every scalar $z$, the morphism $zL_c$ is
anti-Real if and only if $\bar z=-z$.  In particular, the morphism
\[
                         \mathcal K_c=iL_c            \tag{4.23}
\]
is anti-Real.  At $c=1$, $\mathcal K_1^2=-\id$ in every cyclic degree.
\end{proposition}

\begin{proof}
An explicit nonzero degree-zero chain is furnished by the first matrix unit in (4.21).  Applying
translation by $c^{-1}$ recovers it from its $L_c$ image, so $L_c$ is nonzero.  Lemma
\ref{lem:parityintersection} excludes anti-Real parity.  If $zL_c$ is anti-Real, evaluate its
parity identity on a Real conjugate of a nonzero chain.  Real-evenness of $L_c$ reduces the identity
to $\bar zL_cx=-zL_cx$, and nonvanishing gives $\bar z=-z$.  Conversely this scalar identity and
conjugate linearity of $\sigma$ make $zL_c$ anti-Real.  Taking $z=i$ proves (4.23).  Finally
$L_1=\id$ and $i^2=-1$.
\end{proof}

\begin{proof}[Proof of Theorem~\ref{thm:p2}]
Proposition~\ref{prop:landing} proves the genuine full-adeles landing.  Proposition
\ref{prop:fracture} shows that its degree-zero quotient is a cyclic-linear object, not an algebra
quotient; a direct $K_0$ idempotent lift is therefore ill-typed.  The scalar Real-corona Chern
character has anti-Real parity, whereas the CCM Lefschetz correspondence is Real-even and nonzero.
Proposition~\ref{prop:clifford} gives a canonical scalar parity shift and its square law.
\end{proof}

\begin{remark}
Theorem~\ref{thm:p2} is not an equality between a $K_0$ class and a cyclic morphism.  It proves
that such an equality has no type in the exact CCM cokernel and classifies every scalar cyclic
comparison with the required parity: it is $zL_c$ with $\bar z=-z$.  The positive unit choice is
the Clifford suspension $iL_c$.
\end{remark}

\section{Endpoint-holomorphic smoothness}\label{sec:p3}

Let $Z=Z_+\sqcup Z_-\subset\C$ and let $U$ be a connected set containing $Z$ and the relevant
compactifying endpoints.  Define
\[
 \cO_U(Z)=\{f:Z\to\C:\exists h\text{ holomorphic on a neighborhood of }U,
                         \ h|_Z=f\}.                  \tag{5.1}
\]
It is a subalgebra of the full function algebra $\C^Z$.  Give $\C^Z$ the product topology.

\begin{lemma}[Density]\label{lem:density}
$\cO_U(Z)$ is dense in $\C^Z$.
\end{lemma}

\begin{proof}
A basic product neighborhood specifies conditions at finitely many points.  Lagrange interpolation
produces a polynomial agreeing with the prescribed function at those points.  Its restriction lies
in (5.1) and meets the neighborhood exactly.
\end{proof}

Let $e_Z$ be the characteristic function of $Z_+$.  It is an idempotent in $\C^Z$.  Suppose a
holomorphic involution $\rho:U\to U$ preserves $Z$ and exchanges $Z_+$ and $Z_-$.  It induces the
order-two skew crossed products of both $\cO_U(Z)$ and $\C^Z$, and sends $e_Z$ to $1-e_Z$.

\begin{proposition}[Endpoint identity-theorem classification]\label{prop:endpoint}
Assume $Z_+$ and $Z_-$ are nonempty and disjoint.  Suppose that, whenever $Z_\pm$ is infinite, it
accumulates at a point $z_\pm\in U$.  Then
\[
 e_Z\in\cO_U(Z)\quad\Longleftrightarrow\quad
 Z_+\text{ and }Z_-\text{ are finite}.               \tag{5.2}
\]
The equivalent statement holds for the degree-zero element in the reflection crossed product.
\end{proposition}

\begin{proof}
If $Z_+$ is infinite and $h$ represents $e_Z$, then $h-1$ vanishes on a set accumulating inside
the connected holomorphic domain.  The identity theorem gives $h=1$ throughout $U$, contradicting
$h=0$ on nonempty $Z_-$.  The lower-sector argument is identical.  Conversely, for two finite
disjoint sets, Lagrange interpolation gives an entire polynomial taking value one on $Z_+$ and
zero on $Z_-$.  In the skew crossed product, injectivity of the single group coefficient reduces
membership to (5.2).
\end{proof}

We now derive the endpoint hypotheses from completed-function data rather than assuming zero
counting.  Call an entire function $\Lambda$ a \emph{positive-degree self-dual analytic
completion} if there are $C\ge0$ and $0\le a<2$ such that
\begin{align}
 \|\Lambda(z)\|&\le \exp\bigl(C(1+\|z\|)^a\bigr),       \tag{5.3}\\
 \forall A\in\R\quad\exists\sigma\ge1:
 \quad A\sigma&<\log\|\Lambda(\sigma)\|,               \tag{5.4}
\end{align}
every zero lies in $0<\re z<1$, and
\[
             \Lambda(z)=0\Longrightarrow\Lambda(1-z)=0. \tag{5.5}
\]
Condition (5.4) is the only positive-degree input used below.

\begin{lemma}[Finite-zero rigidity]\label{lem:rigidity}
Let $f$ be entire and satisfy \textup{(5.3)}.  If its zero set is finite, then there are a finite
set $T$, multiplicities $m_t$, a constant $c\ne0$, and $B\in\C$ such that
\[
        f(z)=\prod_{t\in T}(z-t)^{m_t}\,c e^{Bz}.     \tag{5.6}
\]
Consequently (5.3) and (5.4) imply that the zero set is infinite.
\end{lemma}

\begin{proof}
Peel the finitely many zeros with their analytic orders.  Outside a ball containing $T$, every
linear factor has norm at least one, so the zero-free quotient inherits (5.3); compactness supplies
the missing bound inside the ball.  An entire logarithm of the quotient has subquadratic real
part, hence is affine, giving (5.6).  Along the positive real ray, the finite product contributes
only logarithmic growth and $e^{Bz}$ only linear logarithmic growth, contradicting (5.4) for a
sufficiently large slope $A$.
\end{proof}

For a completion $\Lambda$, define the centered reciprocal coordinate
\[
               w(\rho)=(\rho-\tfrac12)^{-1}.         \tag{5.7}
\]
Let $Z_+$ and $Z_-$ be images of zeros with positive and negative imaginary part.  Central real
zeros are not placed in either sector; they belong to the fixed channel of Section~\ref{sec:p4}.

\begin{proposition}[Universal compactification]\label{prop:compactification}
For every positive-degree self-dual analytic completion, $Z_+$ and $Z_-$ are infinite, disjoint,
and both accumulate at $0$.  Negation exchanges them.  Therefore $e_Z$ is not smooth in the dense
endpoint-entire coefficient algebra or its order-two reflection crossed product.
\end{proposition}

\begin{proof}
Finite-zero rigidity gives infinitely many zeros.  Entire nontriviality and isolated-zero
theory make the zeros finite in each closed disc.  Strip confinement forces unbounded imaginary
parts.  Reflection (5.5) converts an unbounded sequence of either sign into one of the opposite
sign.  Formula (5.7) tends to zero along both sequences and satisfies
$w(1-\rho)=-w(\rho)$.  Injectivity of reciprocal away from the center proves disjointness.
Proposition~\ref{prop:endpoint} completes the proof.
\end{proof}

\begin{corollary}[Actual completed instances]\label{cor:actual}
The Riemann Xi compactified Real-corona idempotent is not endpoint-holomorphically smooth.  The same
holds for every nontrivial primitive self-dual Dirichlet completion.
\end{corollary}

\begin{proof}
For Xi, infinitude and local finiteness of the complete zero divisor, strip confinement, and the
functional equation instantiate Proposition~\ref{prop:compactification}.  For a nontrivial
primitive Dirichlet character, the completed $L$-function is entire of order at most one and its
positive-real ray growth satisfies (5.4).  Lemma~\ref{lem:rigidity} supplies infinitely many zeros.
The Euler zero-free half-plane and the functional equation confine them to the open unit strip;
self-duality gives (5.5).  Isolated-zero theory supplies local finiteness.
\end{proof}

\begin{proof}[Proof of Theorem~\ref{thm:p3}]
The abstract classification is Proposition~\ref{prop:endpoint}.  Proposition
\ref{prop:compactification} proves the positive-degree statement without assuming zero infinitude,
and Corollary~\ref{cor:actual} gives the actual completed examples.
\end{proof}

\section{Fixed sectors and analytic rank}\label{sec:p4}

Let $T$ be a parameter set.  A completed self-dual analytic family consists of entire nonzero
functions $F_t$, a real center $c$, a nonzero root number $\omega_t$, and the functional equation
\[
                    F_t(2c-z)=\omega_tF_t(z).        \tag{6.1}
\]
Fix $R\ge0$ and count real zeros in $[c-R,c+R]$ with analytic multiplicity.  Isolated-zero theory
makes this a finite multiplicity function $q_t:\R\to_0\N$.  Equation (6.1) preserves analytic
orders, so
\[
                    q_t(2c-x)=q_t(x).                \tag{6.2}
\]
Define
\[
 d(t)=\sum_xq_t(x),\qquad r(t)=q_t(c),\qquad
 u(t)=\sum_{x>c}q_t(x).                              \tag{6.3}
\]

\begin{proposition}[Pointwise defect decomposition]\label{prop:defect}
For every $t$,
\[
                       d(t)=r(t)+2u(t).              \tag{6.4}
\]
Moreover $(-1)^{r(t)}=\omega_t$.
\end{proposition}

\begin{proof}
Reflection (6.2) pairs every $x<c$ with one $x'>c$ of equal multiplicity; the center is the only
fixed point.  This gives (6.4).  Applying (6.1) twice gives $\omega_t^2=1$.  Comparing the first
nonzero Taylor coefficient at $c$ gives $(-1)^{r(t)}=\omega_t$.
\end{proof}

The coordinate $K_0$ model of a finite fixed divisor is the Grothendieck group of finite-support
coordinate idempotents.  The idempotent attached to $q_t$ has value one at $x$ precisely when
$q_t(x)>0$.  Thus its rank is support cardinality, not total multiplicity.

\begin{proposition}[Support--multiplicity fracture]\label{prop:support}
Let $k(t)$ be the central coordinate evaluation of the fixed idempotent.  Then
\[
                       k(t)=\mathbf1_{r(t)>0}.        \tag{6.5}
\]
If $\mu(n)=n-\mathbf1_{n>0}$, then
\[
                       r(t)=k(t)+\mu(r(t)).           \tag{6.6}
\]
The total cyclic fixed-sector rank is the ordinary coordinate $K_0$ rank plus the total
multiplicity excess.
\end{proposition}

\begin{proof}
Coordinate evaluation of an idempotent has rank zero or one according as that coordinate is absent
or present.  This is (6.5); (6.6) is the defining decomposition of a natural number into support
and excess.  Summing it over the finite support gives the final statement.
\end{proof}

\begin{proof}[Proof of Theorem~\ref{thm:p4}]
Subtract (6.4) at $t_0$ from (6.4) at $t_1$ to obtain (1.10).  Subtract (6.6) at the two parameters
and use (6.5) to obtain (1.11).
\end{proof}

The result applies to actual nontrivial primitive self-dual Dirichlet completions: their analytic
orders, root numbers, and real zero packets instantiate (6.1)--(6.3).  No constancy of analytic
rank and no simplicity of central zeros is assumed.

\section{One law, four shadows}\label{sec:synthesis}

We summarize the common mechanism.  Three orthogonal idempotent sectors
\[
              1=e_++e_{\mathrm{fix}}+e_-            \tag{7.1}
\]
are exchanged by the Real involution on $e_\pm$ and fixed on $e_{\mathrm{fix}}$.  The orientation
$m_0=e_+-e_-$ obeys
\[
       m_0^2=1-e_{\mathrm{fix}},\qquad \sigma(m_0)=-m_0. \tag{7.2}
\]
On the free channel, $m_0$ is Clifford: combined with conjugation it squares to $-1$.  On the fixed
channel, (7.2) is the exact square defect.

Theorem~\ref{thm:conservation} supplies the analytic conservation principle.  P1 is its local
extreme: the complete spectral energy is entirely even, so $m_0$ is smooth but its commutator jets
vanish.  P2 is its categorical extreme: after full-adeles descent, the relevant object is a cyclic
module rather than an algebra; Real parity survives only after multiplication by $i$.  P3 is the
holomorphic extreme: an infinite free-orbit sign cannot extend across its accumulation endpoint.
P4 is the fixed collateral: support survives in ordinary $K_0$, while multiplicity survives in the
cyclic rank.

This can be expressed as a conservation table.
\begin{center}
\small
\begin{tabular}{@{}>{\raggedright\arraybackslash}p{0.24\linewidth}
  >{\raggedright\arraybackslash}p{0.25\linewidth}
  >{\raggedright\arraybackslash}p{0.42\linewidth}@{}}
\toprule
Channel & Visible invariant & Forced loss \\
\midrule
Unbounded even & Global spectrum & First local commutator \\
Bounded odd & Local cyclic parity & Unbounded detection \\
Free crossed-product orbit & $iL_c$ & Ordinary trace and direct $K_0$ lift \\
Fixed orbit & Cyclic multiplicity & Multiplicity in coordinate $K_0$ \\
\bottomrule
\end{tabular}
\end{center}
No third parity sector exists because conjugation by an involution has exactly the $\pm1$
eigensplitting (1.1).  The fixed idempotent is not a third operator parity; it is the failure of
the free orientation to square to the unit.

\section{Formal verification}\label{sec:formal}

All introduced mathematical results are formalized in Lean~4 against the repository's pinned
mathlib toolchain.  The facade module is
\begin{center}
\texttt{GaloisForLFunctions.RHTower.RealCoronaFourProblemsResolution}.
\end{center}
It imports the complete dependency graph and exposes stable declarations for Theorems
\ref{thm:conservation}--\ref{thm:p4}.  The formal development includes:
\begin{enumerate}[leftmargin=2.2em]
\item partially defined operators on exact invariant domains, maximal self-adjoint diagonal
operators, and all-orders smooth elements;
\item actual Riemann zero copies with analytic multiplicity and functional-equation reflection;
\item the rational finite-adele topology, normalized Haar character, Bruhat crossed product,
unitization, endpoint cyclic morphisms, Schatten-$1$ Morita fields, strong class-space cyclic
target, and all-degree landing;
\item the honest order-two reflection crossed product of the dense endpoint-holomorphic algebra;
\item analytic orders and finite real-zero packets for self-dual entire families and actual
primitive Dirichlet completions.
\end{enumerate}

No declaration in the paper dependency cone uses \texttt{sorry}, \texttt{admit}, a new axiom,
\texttt{opaque}, or \texttt{unsafe}.  The accompanying \texttt{CLAIM\_TRACE.md} maps every named
result and every substantive displayed identity to exact Lean declarations.  Reproduction uses
targeted commands only:
\begingroup
\scriptsize
\begin{verbatim}
cd formal
lake env lean GaloisForLFunctions/RHTower/RealCoronaFourProblemsResolution.lean
lake build GaloisForLFunctions.RHTower.RealCoronaFourProblemsResolution
python3 scripts/formal_tools.py lint
\end{verbatim}
\endgroup
The final build and lint outputs are recorded in the hostile-review artifact.

\section{Limitations and further structure}

The negative conclusions are classification results for specified natural structures, not
universal impossibility claims about all conceivable models.  Theorem~\ref{thm:p1} concerns the
canonical complete-divisor diagonal and reflection-compatible local series.  A doubled or
non-diagonal spectral triple may introduce a bounded odd component, but Theorem
\ref{thm:conservation} then identifies precisely where the new local information enters.

Theorem~\ref{thm:p2} does not construct an algebra structure on a cyclic cokernel; it proves that
the representative product cannot descend and classifies the surviving scalar cyclic-module
morphisms.  Passage from such a morphism to a fully developed entire bivariant Chern character
would require a chosen bornological completion and its homotopy theory.  The parity classification
and square law proved here are independent of that choice.

The endpoint algebra of Theorem~\ref{thm:p3} is natural because it simultaneously contains all
polynomial restrictions, is dense for pointwise convergence, is preserved by the holomorphic
functional-equation involution, and tests analytic continuation through the compactifying
endpoint.  Other smooth structures may be studied, but they answer a different regularity
question.  Higher-rank automorphic completions fall under Proposition~\ref{prop:compactification}
once (5.3)--(5.5) are supplied; actual higher-rank automorphic representation theory is not yet
part of the Lean library used here.

Finally, the fixed $K_0$ class cannot recover analytic rank beyond support without enrichment.
Proposition~\ref{prop:support} suggests the minimal enrichment: retain the local primary length, or
equivalently pair the idempotent with the multiplicity-sensitive cyclic trace.  This is exactly the
collateral channel predicted by (7.2).

\begin{thebibliography}{99}

\bibitem{Carboni}
G.~Carboni, J.~A.~Guccione, and J.~J.~Guccione,
\emph{Cyclic homology of crossed products},
J. Algebra \textbf{321} (2009), 2827--2860.

\bibitem{CM}
A.~Connes and H.~Moscovici,
\emph{The local index formula in noncommutative geometry},
Geom. Funct. Anal. \textbf{5} (1995), 174--243.

\bibitem{CCM}
A.~Connes, C.~Consani, and M.~Marcolli,
\emph{The Weil proof and the geometry of the adeles class space},
in \emph{Noncommutative Geometry, Arithmetic, and Related Topics},
Johns Hopkins Univ. Press, 2011, 339--405;
\href{https://arxiv.org/abs/math/0703392}{arXiv:math/0703392}.

\bibitem{CCMZeta}
A.~Connes, C.~Consani, and H.~Moscovici,
\emph{Zeta spectral triples},
preprint (2025),
\href{https://arxiv.org/abs/2511.22755}{arXiv:2511.22755}.

\bibitem{CuntzQuillen}
J.~Cuntz and D.~Quillen,
\emph{Excision in bivariant periodic cyclic cohomology},
Invent. Math. \textbf{127} (1997), 67--98.

\bibitem{JLO}
A.~Jaffe, A.~Lesniewski, and K.~Osterwalder,
\emph{Quantum $K$-theory. I. The Chern character},
Comm. Math. Phys. \textbf{118} (1988), 1--14.

\bibitem{Meyer}
R.~Meyer,
\emph{On a representation of the idele class group related to primes and zeros of
$L$-functions},
Duke Math. J. \textbf{127} (2005), 519--595.

\bibitem{Moscovici}
H.~Moscovici,
\emph{Local index formula and twisted spectral triples},
Quanta of Maths, Clay Math. Proc. \textbf{11}, Amer. Math. Soc., 2010, 465--500;
\href{https://arxiv.org/abs/0902.0835}{arXiv:0902.0835}.

\bibitem{Paugam}
F.~Paugam,
\emph{Sym\'etries spectrales des fonctions z\^etas},
J. Th\'eor. Nombres Bordeaux \textbf{21} (2009), 713--720.

\end{thebibliography}

\end{document}
